\documentclass{article}
\usepackage[T2A]{fontenc}
\usepackage[utf8]{inputenc}
\usepackage{amsthm}
\usepackage{amsmath}
\usepackage{amssymb}
\usepackage{amsfonts}
\usepackage{mathrsfs}
\usepackage[12pt]{extsizes}
\usepackage{fancyvrb}
\usepackage{indentfirst}
\usepackage[
  left=2cm, right=2cm, top=2cm, bottom=2cm, headsep=0.2cm, footskip=0.6cm, bindingoffset=0cm
]{geometry}
\usepackage[english,russian]{babel}

\begin{document}
\section*{Вариант 15}

После линеаризации и выполнения интегрального преобразования Лапласа по времени \( t \)
\[
f(t) \to f(\lambda) = \int_0^\infty f(t) e^{-\lambda t} \, dt
\]
динамическая модель КДС представляется в виде матрицы передаточных функций \( \Phi(\lambda, \mathbf{p}, \mathbf{s}) \)

\begin{equation}
\tilde{Y}(\lambda) = \Phi(\lambda, \mathbf{p}, \mathbf{s}) \tilde{X}(\lambda), \qquad 
\Phi(\lambda, \mathbf{p}, \mathbf{s}) = \bigl[ \Phi_{kj}(\lambda, \mathbf{p}, \mathbf{s}) \bigr],
\end{equation}

\begin{equation}
\Phi_{kj}(\lambda, \mathbf{p}, \mathbf{s}) = \frac{Q_{kj}(\lambda, \mathbf{p}, \mathbf{s})}{D(\lambda, \mathbf{p}, \mathbf{s})},
\quad k = 1, 2, \dots, N_y,\; j = 1, 2, \dots, N_x,
\end{equation}

где \( D(\lambda, \mathbf{p}, \mathbf{s}) \) — характеристический квазимногочлен, \( Q_{kj}(\lambda, \mathbf{p}, \mathbf{s}) \) — возмущающие квазимногочлены, 
\( \mathbf{p} = (p_1, \dots, p_{N_p})^T \in \mathbb{R}^{N_p} \) — параметры обратных связей, 
\( \mathbf{s} = (s_1, s_2, \dots, s_{N_s})^T \in \Omega_s \subset \mathbb{R}^{N_s} \) — конструктивные параметры, от которых зависят передаточные функции линеаризованной системы.

\end{document}